\begin{table}[htbp]
\centering
\singlespacing
\footnotesize
\setlength{\tabcolsep}{3.2pt}
\renewcommand{\arraystretch}{1.10}
\caption{Macro-regional out-of-sample yield anomaly forecasts ($\bar{\hat{\epsilon}}$, bu/ac) and realized regional anomalies at harvest closure $H^*(s)$ (2011--2025, 135 study counties).}
\label{tab:yearly_macro_yields}
\begin{tabular*}{\textwidth}{@{\extracolsep{\fill}} c c c c c c c c @{}}
\toprule
\textbf{Year} & \textbf{Realized} & \textbf{Naive} & \textbf{ENet} & \textbf{RF} & \textbf{XGB} & \textbf{LSTM} & \textbf{Ens} \\
\midrule
2012 & -4.07 & 0.00 & \textbf{-3.81} & -2.76 & -2.92 & -4.99 & -4.40 \\
2019 & -1.09 & 0.00 & +0.99 & -0.21 & \textbf{-0.25} & -0.23 & +0.38 \\
2013 & -0.80 & 0.00 & -2.46 & \textbf{-1.33} & -2.21 & -2.44 & -2.45 \\
2011 & -0.76 & 0.00 & \textbf{-0.89} & -1.16 & -1.54 & -1.90 & -1.40 \\
2024 & +0.91 & 0.00 & +2.86 & +1.62 & +2.33 & \textbf{+0.88} & +1.87 \\
2025 & +2.49 & 0.00 & -0.15 & \textbf{+1.48} & +0.62 & +1.15 & +0.50 \\
2014 & +2.89 & 0.00 & -0.56 & +1.18 & \textbf{+1.65} & +0.70 & +0.07 \\
2017 & +3.05 & 0.00 & \textbf{+3.22} & +0.95 & +2.01 & +1.26 & +2.24 \\
2023 & +3.25 & 0.00 & \textbf{+1.64} & +1.04 & +1.60 & +0.59 & +1.12 \\
2022 & +3.48 & 0.00 & +0.19 & +0.80 & +1.28 & \textbf{+2.12} & +1.16 \\
2020 & +3.60 & 0.00 & +0.83 & -0.14 & +0.84 & \textbf{+1.76} & +1.30 \\
2015 & +3.95 & 0.00 & +1.85 & +1.40 & \textbf{+1.98} & +1.32 & +1.58 \\
2018 & +5.25 & 0.00 & \textbf{+2.76} & +1.57 & +2.55 & +2.59 & +2.68 \\
2021 & +6.68 & 0.00 & \textbf{+2.85} & +1.06 & +1.92 & +2.07 & +2.46 \\
2016 & +7.34 & 0.00 & \textbf{+4.09} & +1.69 & +2.88 & +2.69 & +3.39 \\
\midrule
\textbf{RMSE (bu/ac)} & --- & 3.84 & 2.40 & 2.84 & 2.33 & 2.37 & \textbf{2.31} \\
\textbf{MAE (bu/ac)}  & --- & 3.31 & 2.11 & 2.32 & 2.02 & 2.02 & \textbf{2.04} \\
\textbf{Bias (bu/ac)} & --- & -2.41 & \textbf{-1.52} & -1.93 & -1.56 & -1.91 & -1.71 \\
\textbf{Pearson $r$}   & --- & ---  & 0.783 & 0.810 & 0.862 & \textbf{0.909} & 0.881 \\
\bottomrule
\end{tabular*}
\vspace{0.2em}
{\footnotesize \textit{Source: Author's calculation on USDA NASS and ECMWF ERA5-Land data, 2011--2025 out-of-sample test evaluations ($135$ county evaluations per year). Values represent regional arithmetic mean yield anomaly ($\bar{\epsilon} = \frac{1}{N}\sum_i \epsilon_{it}$, bu/ac) across the 135 study counties, ordered by realized anomaly. All machine learning models evaluate the integrated specification (+Lag $\epsilon_{t-1}$). Bold indicates the point forecast closest to realized regional yield anomaly. Naive: deterministic zero-anomaly baseline ($\hat{\epsilon} = 0$, predicting normal weather); ENet: ElasticNet; RF: Random Forest; XGB: XGBoost; LSTM: sequential LSTM with 10-day decad inputs; Ens: 50/50 convex blend of ElasticNet (+Lag) and LSTM (+Lag). At harvest closure $H^*(s)$, evaluations for Minnesota implement the September 30 cutoff ($H=2$) in accordance with Section~\ref{sec:meth_timeline_horizons}. Bottom block summarizes macro-level accuracy metrics against realized regional yield anomalies ($\text{Bias} = \bar{\hat{\epsilon}} - \bar{\epsilon}$).}}
\end{table}
